% =============================================================================
% Rapport de projet — Prédiction du trafic par T-GCN
% Simulation SUMO + OpenStreetMap + MongoDB (Casablanca)
% =============================================================================
% Compilation :
%   pdflatex rapport.tex
%   pdflatex rapport.tex   (2e passe pour la table des matières)
%
% Placez vos captures d'écran dans le dossier rapport/figures/
% =============================================================================

\documentclass[12pt,a4paper]{report}

\usepackage[utf8]{inputenc}
\usepackage[T1]{fontenc}
\usepackage[french]{babel}
\usepackage{geometry}
\usepackage{graphicx}
\usepackage{hyperref}
\usepackage{booktabs}
\usepackage{float}
\usepackage{listings}
\usepackage{xcolor}
\usepackage{amsmath}
\usepackage{caption}
\usepackage{subcaption}
\usepackage{tikz}
\usetikzlibrary{shapes, arrows.meta, positioning}

\geometry{margin=2.5cm}
\hypersetup{colorlinks=true, linkcolor=blue, urlcolor=blue, citecolor=blue}

\lstset{
    basicstyle=\ttfamily\small,
    breaklines=true,
    frame=single,
    backgroundcolor=\color{gray!8},
    keywordstyle=\color{blue},
    commentstyle=\color{green!50!black},
}

\graphicspath{{figures/}}

% --- Informations à compléter ---
\title{%
  \textbf{Prédiction du trafic routier}\\[0.4cm]
  \large Réseau T-GCN, simulation SUMO/OpenStreetMap\\[0.2cm]
  \large et stockage MongoDB — Casablanca%
}
\author{%
  Noms des étudiants\\[0.2cm]
  Encadré par : Nom du professeur\\[0.2cm]
  Établissement — Année universitaire 2025--2026%
}
\date{\today}

\begin{document}

% =============================================================================
\maketitle
\thispagestyle{empty}

\begin{abstract}
Ce projet met en œuvre un pipeline complet de prévision du trafic basé sur le
modèle \textbf{T-GCN} (Temporal Graph Convolutional Network). Dans un premier
temps, le modèle est entraîné sur le jeu de données de référence \textbf{PEMS07}
(région de New York, 883 capteurs). Dans un second temps, une simulation
microscopique du trafic est réalisée sur \textbf{cinq boulevards de Casablanca}
à l'aide de \textbf{SUMO} et de données \textbf{OpenStreetMap}. Les séries
temporelles produites sont stockées dans \textbf{MongoDB}, partitionnées en
ensembles \emph{test} et \emph{T-GCN}, puis exportées pour l'entraînement du
modèle. Le rapport décrit l'architecture, les choix techniques et les résultats
obtenus.
\end{abstract}

\tableofcontents
\newpage

% =============================================================================
\chapter{Introduction}
% =============================================================================

\section{Contexte et objectifs}

La gestion du trafic urbain nécessite des outils capables de \textbf{prédire}
la vitesse, le débit et l'état de congestion sur un réseau routier. Les
réseaux de neurones convolutifs sur graphes temporels (T-GCN) exploitent à la
fois la \textbf{structure spatiale} du réseau (graphe des capteurs ou des
boulevards) et la \textbf{dynamique temporelle} des vitesses mesurées.

Les objectifs du projet sont les suivants :
\begin{itemize}
  \item Implémenter et entraîner un modèle T-GCN sur le benchmark PEMS07 ;
  \item Construire un graphe de distance et un graphe de corrélation du trafic ;
  \item Simuler le trafic sur 5 boulevards parallèles/perpendiculaires à
        Casablanca avec SUMO et OpenStreetMap ;
  \item Stocker les données simulées dans MongoDB avec une séparation
        \emph{test} / \emph{entraînement T-GCN} ;
  \item Calculer les métriques par boulevard (vitesse moyenne, débit, densité,
        état du trafic) ;
  \item Exporter les données au format compatible T-GCN et entraîner le modèle
        sur Casablanca.
\end{itemize}

% -----------------------------------------------------------------------------
% IMAGE 1 : Photo ou carte de Casablanca / contexte du projet
% Fichier suggéré : figures/intro_casablanca.png
% Contenu : carte Google Maps ou OSM du centre de Casablanca avec les 5 boulevards surlignés
% -----------------------------------------------------------------------------
\begin{figure}[H]
  \centering
  \includegraphics[width=0.85\textwidth]{intro_casablanca.png}
  \caption{%
    \textbf{[À insérer]} Carte du centre de Casablanca montrant les cinq boulevards
    étudiés : Mohammed V, Zerktouni, Hassan II, Avenue des FAR et Brahim Roudani.
    Source : OpenStreetMap ou capture Google Maps avec annotations.
  }
  \label{fig:intro-casablanca}
\end{figure}

\section{Organisation du rapport}

Le chapitre~\ref{ch:architecture} présente l'architecture globale.
Le chapitre~\ref{ch:tgcn} détaille le modèle T-GCN et l'expérience PEMS07.
Le chapitre~\ref{ch:simulation} décrit la simulation SUMO/OSM.
Le chapitre~\ref{ch:mongodb} traite du stockage MongoDB.
Le chapitre~\ref{ch:resultats} présente les résultats et visualisations.
Le chapitre~\ref{ch:conclusion} conclut et propose des perspectives.

% =============================================================================
\chapter{Architecture globale du système}
\label{ch:architecture}
% =============================================================================

\section{Vue d'ensemble}

Le projet se compose de deux branches complémentaires :
\begin{enumerate}
  \item \textbf{Branche PEMS07} : données historiques réelles (New York) —
        référence méthodologique et comparaison des graphes distance/corrélation.
  \item \textbf{Branche Casablanca} : données simulées (SUMO + OSM) —
        application locale sur 5 boulevards avec stockage MongoDB.
\end{enumerate}

% Schéma TikZ intégré (pas besoin d'image externe)
\begin{figure}[H]
  \centering
  \begin{tikzpicture}[
      node distance=1.1cm and 1.4cm,
      box/.style={draw, rounded corners, minimum width=2.8cm, minimum height=0.9cm,
                  align=center, font=\small},
      arrow/.style={-{Stealth[length=2.5mm]}, thick}
    ]
    \node[box, fill=blue!10] (osm) {OpenStreetMap\\(Overpass API)};
    \node[box, fill=orange!15, below=of osm] (sumo) {SUMO\\netconvert + TraCI};
    \node[box, fill=green!10, below=of sumo] (mongo) {MongoDB\\(Docker)};
    \node[box, fill=purple!10, below left=1.2cm and 0.5cm of mongo] (test) {Split test\\40\%};
    \node[box, fill=purple!10, below right=1.2cm and 0.5cm of mongo] (tgcn) {Split T-GCN\\60\%};
    \node[box, fill=red!10, below=1.4cm of mongo] (export) {CASABLANCA05\\.npz + .csv};
    \node[box, fill=yellow!15, below=of export] (model) {Modèle T-GCN\\(PyTorch)};

    \node[box, fill=gray!10, right=3.5cm of sumo] (pems) {PEMS07\\883 capteurs};
    \node[box, fill=yellow!15, below=of pems] (model2) {Modèle T-GCN\\(référence)};

    \draw[arrow] (osm) -- (sumo);
    \draw[arrow] (sumo) -- (mongo);
    \draw[arrow] (mongo) -- (test);
    \draw[arrow] (mongo) -- (tgcn);
    \draw[arrow] (tgcn) -- (export);
    \draw[arrow] (export) -- (model);
    \draw[arrow] (pems) -- (model2);
    \draw[dashed, arrow] (model2) -- node[right, font=\scriptsize]{même architecture} (model);
  \end{tikzpicture}
  \caption{Architecture globale du pipeline : simulation Casablanca (gauche) et
           benchmark PEMS07 (droite).}
  \label{fig:architecture}
\end{figure}

% -----------------------------------------------------------------------------
% IMAGE 2 : Capture du dépôt / structure des dossiers dans l'IDE
% Fichier suggéré : figures/structure_projet.png
% Contenu : arborescence TGCN/ (models/, simulation/, data/, train.py...)
% -----------------------------------------------------------------------------
\begin{figure}[H]
  \centering
  \includegraphics[width=0.75\textwidth]{structure_projet.png}
  \caption{%
    \textbf{[À insérer]} Structure du projet dans l'éditeur (VS Code / Cursor) :
    dossiers \texttt{models/}, \texttt{simulation/}, \texttt{utils/}, \texttt{data/}.
  }
  \label{fig:structure}
\end{figure}

\section{Stack technique}

\begin{table}[H]
  \centering
  \caption{Technologies utilisées}
  \label{tab:stack}
  \begin{tabular}{lll}
    \toprule
    Composant        & Technologie        & Rôle \\
    \midrule
    Modèle           & PyTorch            & Implémentation T-GCN \\
    Données NY       & PEMS07 (.npz/.csv) & Benchmark 883 capteurs \\
    Simulation       & Eclipse SUMO 1.27  & Microsimulation trafic \\
    Cartographie     & OpenStreetMap      & Réseau routier Casablanca \\
    Base de données  & MongoDB 7 (Docker) & Stockage séries temporelles \\
    Langage          & Python 3.11        & Pipeline complet \\
    \bottomrule
  \end{tabular}
\end{table}

% =============================================================================
\chapter{Modèle T-GCN et expérience PEMS07}
\label{ch:tgcn}
% =============================================================================

\section{Principe du T-GCN}

Le T-GCN (Temporal Graph Convolutional Network) combine :
\begin{itemize}
  \item un \textbf{GCN} (Graph Convolutional Network) pour la diffusion spatiale
        de l'information entre nœuds du graphe ;
  \item une cellule récurrente de type \textbf{GRU} pour modéliser la dépendance
        temporelle des vitesses.
\end{itemize}

Pour chaque pas de temps $t$, l'entrée est une matrice de vitesses
$X_t \in \mathbb{R}^{N \times 1}$ sur $N$ nœuds (capteurs ou boulevards).
Le modèle prédit les $H$ pas futurs à partir d'une fenêtre de $T_{in}$ pas
historiques.

\subsection{Hyperparamètres}

\begin{table}[H]
  \centering
  \caption{Hyperparamètres principaux (\texttt{config.py})}
  \label{tab:hyperparams}
  \begin{tabular}{ll}
    \toprule
    Paramètre            & Valeur \\
    \midrule
    \texttt{NUM\_NODES}  & 883 (PEMS07) / 5 (Casablanca) \\
    \texttt{SEQ\_LEN}    & 12 pas d'entrée \\
    \texttt{PRED\_LEN}   & 12 pas de prédiction \\
    \texttt{HIDDEN\_DIM} & 64 \\
    \texttt{BATCH\_SIZE} & 64 \\
    \texttt{LEARNING\_RATE} & $10^{-3}$ \\
    Split temporel       & 60\% train / 20\% val / 20\% test \\
    \bottomrule
  \end{tabular}
\end{table}

\section{Construction des graphes}

Deux types de graphes d'adjacence sont implémentés (\texttt{utils/graph\_utils.py}) :

\subsection{Niveau 1 — Graphe de distance}

Les arêtes proviennent du fichier \texttt{PEMS07.csv} (distances entre capteurs).
Le poids est calculé par noyau gaussien :
\begin{equation}
  W_{ij} = \exp\left(-\frac{d_{ij}^2}{\sigma^2}\right)
\end{equation}
puis sparsifié (seuil 0{,}1) et normalisé symétriquement.

\subsection{Niveau 2 — Graphe de corrélation}

Les arêtes sont construites à partir de la corrélation de Pearson des vitesses
sur la partition d'entraînement uniquement (seuil $|r| \geq 0{,}5$).

% -----------------------------------------------------------------------------
% IMAGE 3 : Heatmap du graphe ou comparaison distance vs corrélation
% Fichier suggéré : figures/graphe_adjacency.png
% Contenu : outputs/plots/ généré par evaluate.py ou explain.py
% -----------------------------------------------------------------------------
\begin{figure}[H]
  \centering
  \includegraphics[width=0.7\textwidth]{graphe_adjacency.png}
  \caption{%
    \textbf{[À insérer]} Visualisation du graphe d'adjacence (distance ou
    corrélation). Générer avec \texttt{python explain.py} ou
    \texttt{python evaluate.py --graph correlation}.
  }
  \label{fig:graphe}
\end{figure}

\section{Entraînement et évaluation}

\subsection{Commandes}

\begin{lstlisting}[language=bash]
# Entrainement PEMS07 — graphe distance
python train.py --graph distance

# Entrainement PEMS07 — graphe correlation
python train.py --graph correlation

# Entrainement Casablanca (apres export SUMO)
python train.py --dataset casablanca --graph distance
\end{lstlisting}

\subsection{Métriques d'évaluation}

\begin{itemize}
  \item \textbf{MAE}, \textbf{RMSE}, \textbf{MAPE} (régression) ;
  \item Classification de l'état du trafic : Congested ($\leq 30$ km/h),
        Moderate ($\leq 60$ km/h), Free Flow ($> 60$ km/h).
\end{itemize}

% -----------------------------------------------------------------------------
% IMAGE 4 : Courbe de loss train/val
% Fichier suggéré : figures/loss_curve.png
% Contenu : copier outputs/plots/loss_curve.png
% -----------------------------------------------------------------------------
\begin{figure}[H]
  \centering
  \includegraphics[width=0.75\textwidth]{loss_curve.png}
  \caption{%
    \textbf{[À insérer]} Courbes de perte MAE en entraînement et validation
    (\texttt{outputs/plots/loss\_curve.png}).
  }
  \label{fig:loss}
\end{figure}

% -----------------------------------------------------------------------------
% IMAGE 5 : Prédiction vs vérité terrain
% Fichier suggéré : figures/pred_vs_true.png
% Contenu : outputs/plots/ de evaluate.py
% -----------------------------------------------------------------------------
\begin{figure}[H]
  \centering
  \includegraphics[width=0.85\textwidth]{pred_vs_true.png}
  \caption{%
    \textbf{[À insérer]} Comparaison vitesse prédite vs réelle sur quelques
    capteurs/boulevards (\texttt{outputs/plots/}).
  }
  \label{fig:pred}
\end{figure}

\section{Explicabilité (XAI)}

Le script \texttt{explain.py} réalise une analyse par \textbf{ablation des
voisins} : on masque successivement chaque nœud voisin dans le graphe pour
mesurer son influence sur la prédiction d'un capteur cible.

% -----------------------------------------------------------------------------
% IMAGE 6 : Heatmap XAI / influence des voisins
% Fichier suggéré : figures/xai_ablation.png
% Contenu : outputs/plots/ de explain.py
% -----------------------------------------------------------------------------
\begin{figure}[H]
  \centering
  \includegraphics[width=0.8\textwidth]{xai_ablation.png}
  \caption{%
    \textbf{[À insérer]} Analyse XAI : influence des nœuds voisins sur la
    prédiction (généré par \texttt{python explain.py}).
  }
  \label{fig:xai}
\end{figure}

% =============================================================================
\chapter{Simulation SUMO et OpenStreetMap — Casablanca}
\label{ch:simulation}
% =============================================================================

\section{Zone d'étude et boulevards}

La simulation porte sur une zone de $\approx 3 \times 3$ km au centre de
Casablanca, téléchargée via l'API Overpass (OpenStreetMap).

\begin{table}[H]
  \centering
  \caption{Cinq boulevards modélisés}
  \label{tab:boulevards}
  \begin{tabular}{clcc}
    \toprule
    ID & Boulevard              & Orientation  & Arêtes SUMO \\
    \midrule
    0  & Mohammed V             & Nord--Sud    & $\sim$185 \\
    1  & Zerktouni              & Nord--Sud    & $\sim$168 \\
    2  & Hassan II              & Est--Ouest   & $\sim$162 \\
    3  & Avenue des FAR         & Est--Ouest   & $\sim$162 \\
    4  & Brahim Roudani         & Connecteur   & $\sim$162 \\
    \bottomrule
  \end{tabular}
\end{table}

\section{Pipeline de simulation}

\begin{enumerate}
  \item \textbf{Téléchargement OSM} (\texttt{simulation/osm/download\_osm.py}) :
        requête Overpass sur la bbox Casablanca ;
  \item \textbf{Construction réseau} (\texttt{netconvert}) : conversion
        OSM $\rightarrow$ \texttt{casablanca.net.xml} ;
  \item \textbf{Demande de trafic} (\texttt{randomTrips.py}) : génération de
        $\sim$2500 véhicules sur 3600 s simulées ;
  \item \textbf{Détecteurs E1} : boucles inductive sur chaque boulevard ;
  \item \textbf{Simulation TraCI} : collecte vitesse, débit, densité toutes les
        60 s (après 300 s de chauffe).
\end{enumerate}

% -----------------------------------------------------------------------------
% IMAGE 7 : Terminal — téléchargement OSM réussi
% Fichier suggéré : figures/osm_download.png
% Contenu : capture "Fichier sauvegardé: casablanca.osm.xml (675.4 KB)"
% -----------------------------------------------------------------------------
\begin{figure}[H]
  \centering
  \includegraphics[width=0.9\textwidth]{osm_download.png}
  \caption{%
    \textbf{[À insérer]} Capture terminal : téléchargement réussi du fichier OSM
    (\texttt{python -m simulation.osm.download\_osm}).
  }
  \label{fig:osm}
\end{figure}

% -----------------------------------------------------------------------------
% IMAGE 8 : SUMO-GUI — réseau routier Casablanca
% Fichier suggéré : figures/sumo_reseau.png
% Contenu : capture SUMO-GUI montrant le réseau (vue dézoomée ou zoomée)
% -----------------------------------------------------------------------------
\begin{figure}[H]
  \centering
  \includegraphics[width=0.85\textwidth]{sumo_reseau.png}
  \caption{%
    \textbf{[À insérer]} Réseau routier SUMO généré depuis OpenStreetMap
    (capture SUMO-GUI ou \texttt{casablanca.net.xml} visualisé).
  }
  \label{fig:sumo-reseau}
\end{figure}

% -----------------------------------------------------------------------------
% IMAGE 9 : SUMO-GUI avec véhicules (zoom centre-ville)
% Fichier suggéré : figures/sumo_voitures.png
% Contenu : SUMO-GUI zoomé sur Mohammed V avec voitures visibles, t > 300s
% -----------------------------------------------------------------------------
\begin{figure}[H]
  \centering
  \includegraphics[width=0.85\textwidth]{sumo_voitures.png}
  \caption{%
    \textbf{[À insérer]} Simulation SUMO avec trafic visible (zoom sur le
    centre-ville, temps simulé $t > 300$ s).
    Commande : \texttt{python -m simulation.run\_pipeline --step simulate --gui}.
  }
  \label{fig:sumo-voitures}
\end{figure}

% -----------------------------------------------------------------------------
% IMAGE 10 : Terminal — mapping boulevards → arêtes
% Fichier suggéré : figures/sumo_mapping.png
% Contenu : sortie "[Boulevard Mohammed V] → 185 arêtes" etc.
% -----------------------------------------------------------------------------
\begin{figure}[H]
  \centering
  \includegraphics[width=0.9\textwidth]{sumo_mapping.png}
  \caption{%
    \textbf{[À insérer]} Association automatique des boulevards aux arêtes SUMO
    par correspondance de noms OSM.
  }
  \label{fig:mapping}
\end{figure}

\section{Paramètres de simulation}

\begin{table}[H]
  \centering
  \caption{Paramètres SUMO (\texttt{simulation/config.py})}
  \label{tab:sumo}
  \begin{tabular}{ll}
    \toprule
    Paramètre               & Valeur \\
    \midrule
    Durée simulée           & 3600 s (1 h) \\
    Pas d'échantillonnage   & 60 s \\
    Période de chauffe      & 300 s \\
    Véhicules cibles        & 2500 \\
    \bottomrule
  \end{tabular}
\end{table}

% =============================================================================
\chapter{Stockage MongoDB}
\label{ch:mongodb}
% =============================================================================

\section{Architecture de la base}

MongoDB est déployé via Docker :
\begin{lstlisting}[language=bash]
docker run -d --name mongo-tgcn -p 27017:27017 mongo:7
\end{lstlisting}

Base de données : \texttt{tgcn\_casablanca}.

\begin{table}[H]
  \centering
  \caption{Collections MongoDB}
  \label{tab:mongo}
  \begin{tabular}{ll}
    \toprule
    Collection           & Contenu \\
    \midrule
    \texttt{boulevards}  & Définition des 5 boulevards \\
    \texttt{simulation\_runs} & Métadonnées de chaque exécution \\
    \texttt{speed\_timeseries} & Vitesse, débit, densité par pas \\
    \texttt{boulevard\_metrics} & Métriques agrégées par boulevard \\
    \texttt{dataset\_splits} & Répartition test / T-GCN \\
    \bottomrule
  \end{tabular}
\end{table}

\section{Partitionnement des données}

Comme demandé dans le cahier des charges :
\begin{itemize}
  \item \textbf{40\% test} : évaluation, visualisation des métriques ;
  \item \textbf{60\% T-GCN} : export vers \texttt{CASABLANCA05.npz} pour
        l'entraînement.
\end{itemize}

Le split est \textbf{temporel} : les premiers pas de temps vont au test, les
derniers à l'entraînement.

% -----------------------------------------------------------------------------
% IMAGE 11 : MongoDB Compass ou terminal mongo
% Fichier suggéré : figures/mongodb_collections.png
% Contenu : MongoDB Compass montrant les collections et un document exemple
%           OU capture "docker ps" + "Connexion OK"
% -----------------------------------------------------------------------------
\begin{figure}[H]
  \centering
  \includegraphics[width=0.85\textwidth]{mongodb_collections.png}
  \caption{%
    \textbf{[À insérer]} Base MongoDB \texttt{tgcn\_casablanca} : collections
    et exemple de document \texttt{speed\_timeseries}.
    Outil : MongoDB Compass (\url{https://www.mongodb.com/products/compass}).
  }
  \label{fig:mongo}
\end{figure}

% -----------------------------------------------------------------------------
% IMAGE 12 : Terminal — métriques par boulevard
% Fichier suggéré : figures/metrics_terminal.png
% Contenu : sortie de python -m simulation.run_pipeline --metrics
% -----------------------------------------------------------------------------
\begin{figure}[H]
  \centering
  \includegraphics[width=0.85\textwidth]{metrics_terminal.png}
  \caption{%
    \textbf{[À insérer]} Tableau des métriques par boulevard affiché dans le
    terminal (\texttt{python -m simulation.run\_pipeline --metrics}).
  }
  \label{fig:metrics}
\end{figure}

% =============================================================================
\chapter{Résultats et visualisations}
\label{ch:resultats}
% =============================================================================

\section{Métriques de trafic par boulevard}

Pour chaque boulevard et chaque pas de temps, le pipeline calcule :

\begin{table}[H]
  \centering
  \caption{Métriques collectées}
  \label{tab:metriques}
  \begin{tabular}{lll}
    \toprule
    Métrique            & Unité          & Description \\
    \midrule
    Vitesse moyenne     & km/h           & \texttt{avg\_speed\_kmh} \\
    Débit               & véh./h         & \texttt{flow\_vph} \\
    Densité             & véh./km        & \texttt{density\_vpk} \\
    Occupation          & \%             & \texttt{occupancy\_pct} \\
    État du trafic      & ---            & Congested / Moderate / Free Flow \\
    \bottomrule
  \end{tabular}
\end{table}

% -----------------------------------------------------------------------------
% IMAGE 13 : Graphique vitesses par boulevard dans le temps
% Fichier suggéré : figures/speed_timeseries.png
% Contenu : simulation/outputs/plots/speed_test_*.png ou speed_tgcn_*.png
%           généré par python -m simulation.view_metrics
% -----------------------------------------------------------------------------
\begin{figure}[H]
  \centering
  \includegraphics[width=0.95\textwidth]{speed_timeseries.png}
  \caption{%
    \textbf{[À insérer]} Évolution de la vitesse moyenne par boulevard
    (split test ou T-GCN). Générer avec
    \texttt{python -m simulation.view\_metrics}.
  }
  \label{fig:speed-ts}
\end{figure}

\section{Export vers T-GCN}

Le module \texttt{simulation/export\_tgcn.py} produit :
\begin{itemize}
  \item \texttt{data/CASABLANCA05.npz} : matrice $[T, 5]$ des vitesses ;
  \item \texttt{data/CASABLANCA05.csv} : graphe de distances géodésiques
        entre boulevards (format compatible \texttt{graph\_utils.py}).
\end{itemize}

% -----------------------------------------------------------------------------
% IMAGE 14 : Comparaison PEMS07 vs Casablanca (optionnel)
% Fichier suggéré : figures/comparaison_datasets.png
% Contenu : tableau ou graphique comparant N=883 vs N=5, NY vs Casablanca
% -----------------------------------------------------------------------------

\section{Lien PEMS07 $\rightarrow$ Casablanca}

\begin{table}[H]
  \centering
  \caption{Comparaison des deux jeux de données}
  \label{tab:compare}
  \begin{tabular}{lcc}
    \toprule
    Caractéristique  & PEMS07        & Casablanca \\
    \midrule
    Région           & New York      & Casablanca \\
    Nombre de nœuds  & 883           & 5 \\
    Source           & Capteurs réels & Simulation SUMO \\
    Format           & $[T, N]$      & $[T, N]$ \\
    Fenêtre          & 12 / 12       & 12 / 12 \\
    Stockage         & Fichiers locaux & MongoDB + NPZ \\
    \bottomrule
  \end{tabular}
\end{table}

% =============================================================================
\chapter{Conclusion et perspectives}
\label{ch:conclusion}
% =============================================================================

\section{Bilan}

Ce projet a permis de :
\begin{itemize}
  \item implémenter un modèle T-GCN complet (entraînement, évaluation, XAI) ;
  \item comparer deux stratégies de graphe (distance et corrélation) ;
  \item simuler le trafic sur 5 boulevards de Casablanca à partir d'OpenStreetMap ;
  \item stocker et partitionner les données dans MongoDB ;
  \item exporter un jeu de données local compatible avec le pipeline T-GCN.
\end{itemize}

\section{Perspectives}

\begin{itemize}
  \item Augmenter le nombre de nœuds (plus de boulevards ou capteurs par boulevard) ;
  \item Calibrer la demande SUMO sur des données réelles de Casablanca ;
  \item Déployer une API temps réel (FastAPI) connectée à MongoDB ;
  \item Comparer les performances T-GCN PEMS07 vs Casablanca dans un tableau unifié ;
  \item Intégrer la prédiction en boucle fermée dans SUMO (trafic adaptatif).
\end{itemize}

% =============================================================================
\appendix
\chapter{Guide des figures à capturer}
% =============================================================================

Ce tableau récapitule \textbf{toutes les images} à placer dans
\texttt{rapport/figures/}.

\begin{table}[H]
  \centering
  \small
  \begin{tabular}{clp{6.5cm}}
    \toprule
    Fichier & Figure & Comment l'obtenir \\
    \midrule
    \texttt{intro\_casablanca.png}      & Fig.~\ref{fig:intro-casablanca} & Carte OSM/Google Maps annotée \\
    \texttt{structure\_projet.png}      & Fig.~\ref{fig:structure}        & Capture IDE (arborescence) \\
    \texttt{graphe\_adjacency.png}       & Fig.~\ref{fig:graphe}           & \texttt{python explain.py} \\
    \texttt{loss\_curve.png}             & Fig.~\ref{fig:loss}             & \texttt{outputs/plots/loss\_curve.png} \\
    \texttt{pred\_vs\_true.png}           & Fig.~\ref{fig:pred}             & \texttt{python evaluate.py} \\
    \texttt{xai\_ablation.png}           & Fig.~\ref{fig:xai}              & \texttt{python explain.py} \\
    \texttt{osm\_download.png}           & Fig.~\ref{fig:osm}              & Terminal OSM download \\
    \texttt{sumo\_reseau.png}            & Fig.~\ref{fig:sumo-reseau}      & SUMO-GUI réseau vide \\
    \texttt{sumo\_voitures.png}          & Fig.~\ref{fig:sumo-voitures}    & SUMO-GUI avec trafic \\
    \texttt{sumo\_mapping.png}           & Fig.~\ref{fig:mapping}          & Terminal mapping arêtes \\
    \texttt{mongodb\_collections.png}    & Fig.~\ref{fig:mongo}            & MongoDB Compass \\
    \texttt{metrics\_terminal.png}       & Fig.~\ref{fig:metrics}          & \texttt{--metrics} \\
    \texttt{speed\_timeseries.png}       & Fig.~\ref{fig:speed-ts}         & \texttt{view\_metrics} \\
    \bottomrule
  \end{tabular}
  \caption{Liste complète des captures d'écran à insérer}
\end{table}

\chapter{Commandes utiles}

\begin{lstlisting}[language=bash]
# Environnement
pip install -r requirements.txt
docker start mongo-tgcn

# Pipeline Casablanca complet
python -m simulation.run_pipeline

# Etapes separees
python -m simulation.run_pipeline --step osm
python -m simulation.run_pipeline --step sumo
python -m simulation.run_pipeline --step simulate
python -m simulation.run_pipeline --step simulate --gui
python -m simulation.run_pipeline --step export
python -m simulation.run_pipeline --metrics

# T-GCN
python train.py --graph distance
python train.py --dataset casablanca
python evaluate.py --dataset casablanca
python explain.py
python -m simulation.view_metrics
\end{lstlisting}

\end{document}
